% \begin{figure*}
%     \centering
%     \includegraphics[width=0.7\linewidth]{sections/rubricifying_graph1.pdf}
%     \caption{This shows the various comparisons made as explained in Section~\ref{sec:experiments}.}
%     \label{fig:experiment_design}
% \end{figure*}

\section{Experiment Setup}
% \alfredo{This section no longer flows right}
% A prompt is defined as all the text provided to a rater, which may include instructions, examples, additional context, and a rubric (i.e., scoring guidelines). This work explores three components of prompt design and evaluation: agreement level of human raters, decomposition level of a prompt, and generality level of a prompt. 
% \jessica{\textbf{Agreement level} refers to the degree of consensus among human raters, categorized as \textit{full agreement}, \textit{full disagreement}, or in the case of instruction-following, \textit{partial agreement} between the human raters.}

% \subsection{Hypotheses}

Our hypotheses focus on the criteria presentation of rubrics for autoraters:
% (1) Autoraters and humans will exhibit higher alignment on texts with \textbf{full agreement} than on texts with \textbf{partial} or \textbf{full disagreement} (texts with higher inter-human agreement should be more straightforward to evaluate for the autorater based on the given rubric);
(1) \textbf{Edited} prompts will improve autorater alignment with human ratings over the original prompts (humans and autoraters may require different prompts due to the way that they process information); and 
(2) \textbf{Analytic} rubrics will improve autorater alignment with human ratings over \textbf{holistic} rubrics (the decomposition of holistic constructs into discrete components within analytic rubrics is expected to improve autorater alignment by simplifying the evaluation task).

\subsection{Experiments}
\label{sec:experiments}

This work examines four different scores: human ratings on holistic ($\textnormal{H}_{H}$) and analytic prompts ($\textnormal{H}_{A}$), and autorater ratings on holistic ($\textnormal{LLM}_{H}$) and analytic prompts ($\textnormal{LLM}_{A}$). These scores facilitate two comparisons: \textbf{$\Delta$Rater}, which assesses human-autorater agreement when using the same type of rubric, \textbf{$\Delta$Rubric}, which investigates the impact of varying the type rubric while keeping the type of rater constant, revealing how different rubrics influence ratings even when measuring the same criteria. These comparisons are depicted in Figure \ref{fig:experiment_design}. Additional analysis on \textbf{$\Delta$Rater+Rubric}, varying both the type of rubric and type of rater, is explored in Appendix~\ref{sec:other_obs} for completion.

Additionally, this study compares two approaches: presenting the autorater with the original human rubric versus a modified rubric, or \textbf{edited rubric}, designed to enhance agreement. Modifications include adding additional context (if available), incorporating examples, and reducing positional bias. Prior work \citep{mansour-etal-2024-large} studied adding a rubric and then adding an example into the prompt, and found that ChatGPT\footnote{\label{note}\url{https://platform.openai.com/docs/models}} benefited from having rubrics and examples, whereas Llama-2-13b-chat-hf\footnote{\url{https://huggingface.co/meta-llama/Llama-2-13b-chat-hf}} did not for some cases. This work uses several modification conditions. First, the original holistic prompts for the AES task include a large number of examples (10-18) - edited holistic prompts reduces the examples to a representative set of three examples (\textbf{3ex}; high, medium, and low scores). This serves two purposes: (1) it creates a more direct comparison with the \textbf{edited} analytic prompts, which also use three examples, and (2) it avoids an excessively long prompt for the model. Second, the analytic rubric prompt formats of presenting all sub-criteria in a single API call (batch), in separate API calls to mitigate potential positional bias (separate), and a combination of enhancements such as incorporating additional context, the "3ex" set (as the original analytic rubrics had no examples), and the "separate" API call strategy (edited) are tested. Prompts are provided in Appendix~\ref{sec:appendix_B}. 

\section{Methods and Materials}

\subsection{Datasets}
\label{sec:datasets}

Each hypothesis is tested across the domains of AES and IF to determine whether any observed effect is context-dependent. Human annotations are provided by each dataset.

\paragraph{Automatic Essay Scoring (AES).} Automatic essay scoring is a well-established field, with the Automated Student Assessment Prize (ASAP) as a primary dataset \citep{asap-aes}. ASAP contains over 21,000 essays over 8 essay prompts, each with a holistic rubric used by trained human raters. To provide further diversity within AES, this study uses essay prompts 1, 4, and 6, which represent diverse essay prompt types (argumentative, source-dependent, source-dependent) and rubric types (general, general, task-specific) shown in Appendix Table~\ref{tab:asap_prompts}. The essay prompts also provide single holistic scores, and integer scores from 1 to 5 (essay prompt 6's scores are condensed to integers). 

The ASAP++ dataset \citep{mathias-bhattacharyya-2018-asap} provides corresponding analytic rubrics for these essay prompts. Except for the holistic rubric for essay prompt 6, all other rubrics are general. It is important to note while both ASAP and ASAP++ evaluate essay quality, the specific sub-criteria for essay quality differ (ex. ``audience awareness'' is mentioned in ASAP but not ASAP++ and ``conventions'' is mentioned in ASAP++ but ASAP explicitly asks raters to not consider this). This discrepancy may lead to imperfect agreement between ASAP and combined ASAP++ ratings. Additionally, the ASAP++ analytic rubrics do not include examples, so three examples per essay prompt were selected from ASAP++ and excluded from analysis. These examples represent high, medium, and low scores across all analytic sub-criteria (ex. a high scoring essay contains the highest scores for all analytic sub-criteria).


\paragraph{Instruction-following (IF).} 
Compared to AES, 
instruction-following work has decomposed instructions into a series of questions that isolate individual criteria, but with no evaluation differences in the rubric. The InfoBench \citep{qin-etal-2024-infobench} dataset consists of both easy and hard instructions, along with outputs from five popular LLMs. Outputs are holistically scored on a scale of one to five with a general rubric, and analytically scored with binary yes/no responses to decomposed instructions with task-specific rubrics. In the dataset, each output was annotated with both scores by three expert annotators using provided rubrics. Expert annotators annotated 50 instructions with 5 generations each. The same example selection methodology used for AES was used to select input, output, instruction tuples from the remaining dataset that were not annotated by the expert annotators, since neither rubric in IF provides examples. The decomposed automatic instruction annotations released by the InfoBench authors served as a starting point to explore the tuples. The resulting example outputs were generated from three models: GPT-3.5-turbo\footref{note}, Alpaca-7b~\citep{taori2023alpaca}, and GPT-4\footref{note}\citep{openai2024gpt4technicalreport}), and represent high, medium, and low instruction-following ratios (100\%, 50\%, 0\%). This is translated into holistic scores of 5, 3, and 1, respectively, which were validated by the authors.

We select AES and IF for complementary purposes. The AES dataset provides validity through its multiple rubric variations (examples) all measuring the same underlying construct, essay quality, allowing us to observe the different effects of different operationalizations of the same evaluation task. Although using ASAP and ASAP++ simultaneously may result in a noisy comparison, it provides preliminary evidence for the behavior. Conversely the InfoBench dataset offers a direct 1:1 comparison holistic and analytic rubrics that measure identical criteria. Together, these datasets enable us to distinguish between effects that are domain specific versus those that generalize across evaluation contexts/tasks.

\subsection{Autoraters}
Experiments were conducted using gpt-4o-2024-11-20\footref{note}\citep{openai2024gpt4ocard} as the autorater, which has demonstrated high alignment with human evaluations across various tasks \citep{chan:chateval, zhou:sotopia}, and Llama-3.1-70B-Instruct \citep{grattafiori2024llama}. 

% The first hypothesis is because of the inherent issues still present in autorater CoT where autorater reasoning has not reached guaranteed success. By reducing the inference required to make decisions into simpler tasks such as presence of a characteristic or looking at sub-characteristics of a more abstract characteristic, an autorater may have less variable decisions and therefore be more similar to humans. Although analytic rubrics attempt to decompose the task, they may still be high inference, which does not decrease cognitive load.

% \subsection{Autorater Output}
Autorater ratings are calculated using a probability weighting scheme similar to \citet{huynh2023understanding}, given by the equation: $r = \sum_{i=0}^{n}\frac{{p_i}}{\sum_{j=0}^{n} {p_j}} *s_i$, where $r$ is the final rating from the autorater, $n$ represents the number of possible scores given by the rubric, $p_i$ represents the exponential of the log probability score given by the autorater, and $s_i$ represents the integer score outputted by the autorater that corresponds to the log probability. In IF for analytic prompts, $s_i$ is represented by 1 for ``yes'' and 0 for ``no''.

\begin{table*}[!ht]
    \small
    \centering
    \begin{tabular}{|w{c}{0.05cm}|w{c}{0.3cm}||c|c||c|c|c||c|c||c|c|c|}
    \hline
        & & \multicolumn{5}{c||}{GPT-4o} & \multicolumn{5}{c|}{Llama}\\
        \hline
        & & \multicolumn{10}{c|}{$\Delta\textbf{rater}$} \\
        \hline
        & & \multicolumn{2}{c||}{$\textnormal{H}_{H} \to \textnormal{LLM}_H$}  & \multicolumn{3}{c||}{$\textnormal{H}_{A} \to \textnormal{LLM}_A$} & \multicolumn{2}{c||}{$\textnormal{H}_{H} \to \textnormal{LLM}_H$}  & \multicolumn{3}{c|}{$\textnormal{H}_{A} \to \textnormal{LLM}_A$}\\
        \hline
        P & C. & Full & 3ex & Sep. & Bat. & Edited& Full & 3ex & Sep. & Bat. & Edited \\
        \hline
        \multirow{6}{*}{1}& H. & $0.437^{\dagger}$ & $0.387$ & - & - & - & $0.576$ & $0.571$ & - & - & -\\
        \cline{2-12}
        & Ide. & -& -& $0.474^{\star}$ & $0.464$ & $0.552^{s,b\uparrow}$ & -& -& $0.552$ & $0.559$ & $0.552$\\
        \cline{2-12}
        & Or. & -& -& $0.497^{\star}$ & $0.448$ & $0.544^{s,b\uparrow}$ & -& -&$0.547^{\star}$ & $0.524$ & $0.529^{s\downarrow}$\\
        \cline{2-12}
        & WC & -& -& $0.452^{\star}$ & $0.439$ & $0.554^{s,b\uparrow}$ & -& -& $0.547^{\star}$ & $0.527$ & $0.507^{s\downarrow}$\\
        \cline{2-12}
        & SF & -& -& $0.459$ & $0.455$ & $0.553^{s,b\uparrow}$ & -& -& $0.559^{\star}$ & $0.512$ & $0.523^{s\downarrow}$\\
        \cline{2-12}
        & Cv. & -& -& $0.362$ & $0.388^{\star}$ & $0.472^{s,b\uparrow}$ & -& -& $0.491$& $0.477$ & $0.487$\\
        \hline\hline
        \multirow{5}{*}{4}& H. & $0.695^{\dagger}$ & $0.687$ & - & - & - & $0.700$ & $0.699$ & - & - & -\\
        \cline{2-12}
        & Ct. & -& -& $0.673$ & $0.696^{\star}$ & $0.698^{s\uparrow}$& -& -& $0.698^{\star}$ & $0.693$ & $0.701$\\
        \cline{2-12}
        & PA & -& -& $0.662$ & $0.677^{\star}$ & $0.680^{s\uparrow}$& -& -& $0.683^{\star}$ & $0.665$ & $0.682^{b\uparrow}$\\
        \cline{2-12}
        & La. & -& -& $0.627^{\star}$ & $0.595$ & $0.639^{s,b\uparrow}$& -& -& $0.630^{\star}$ & $0.573$ & $0.621^{b\uparrow}$\\
        \cline{2-12}
        & Na. & -& -& $0.669^{\star}$ & $0.652$ & $0.682^{s,b\uparrow}$& -& -& $0.674^{\star}$ & $0.642$ & $0.669^{b\uparrow}$\\
        \hline\hline
        \multirow{6}{*}{6} & H. & $0.629$ & $0.644^{\dagger}$ & - & - & -& $0.666$ & $0.680^{\dagger}$ & - & - & -\\
        \cline{2-12}
        & Ct. & -& -& $0.610$ & $0.652^{\star}$ & $0.676^{s,b\uparrow}$& -& -& $0.680^{\star}$ & $0.648$ & $0.694^{s,b\uparrow}$\\
        \cline{2-12}
        & PA & -& -& $0.605$ & $0.608$ & $0.666^{s,b\uparrow}$ & -& -& $0.652^{\star}$ & $0.601$ & $0.668^{s,b\uparrow}$\\
        \cline{2-12}
        & La. & -& -& $0.524^{\star}$ & $0.510$ & $0.546^{s,b\uparrow}$& -& -& $0.542^{\star}$ & $0.496$ & $0.547^{b\uparrow}$\\
        \cline{2-12}
        & Na. & -& -& $0.562$ & $0.557$ & $0.590^{s,b\uparrow}$& -& -& $0.579^{\star}$ & $0.547$ & $0.589^{b\uparrow}$\\
        \hline\hline
        & & \multicolumn{10}{c|}{$\Delta\textbf{rubric}$} \\
        \hline
        & & \multicolumn{2}{c||}{$\textnormal{H}_{H} \to \textnormal{H}_A$}& \multicolumn{3}{c||}{$\textnormal{LLM}_H \to \textnormal{LLM}_A$} & \multicolumn{2}{c||}{$\textnormal{H}_{H} \to \textnormal{H}_A$}& \multicolumn{3}{c|}{$\textnormal{LLM}_H \to \textnormal{LLM}_A$}\\
        \hline
        P & C. & Full & 3ex & Sep. & Bat. & Edited& Full & 3ex & Sep. & Bat. & Edited \\
        \hline
        \multirow{2}{*}{1} & full &$0.591$ &- & $0.756$& $0.750$& $0.839^{s,b\uparrow}$ &- &- & $0.685^{\star}$ & $0.671$& $0.711^{s,b\uparrow}$\\
        \cline{2-12}
        & 3ex &- &- & $0.792$& $0.789$& $0.838^{s,b\uparrow}$ &- &- & $0.736^{\star}$& $0.713$& $0.768^{s,b\uparrow}$\\
        \hline
        \multirow{2}{*}{4} & full &$0.656$ &- &$0.838^{\star}$ & $0.817$ & $0.876^{s,b\uparrow}$ &- &- & $0.854^{\star}$& $0.780$ & $0.869^{s,b\uparrow}$\\
        \cline{2-12}
        & 3ex&- &- & $0.847^{\star}$& $0.823$& $0.880^{s,b\uparrow}$ &- &- & $0.861^{\star}$& $0.784$& $0.877^{s,b\uparrow}$\\
        \hline
        \multirow{2}{*}{6} & full &$0.681$ &- & $0.781^{\star}$& $0.773$& $0.842^{s,b\uparrow}$ & -&- & $0.803^{\star}$& $0.748$& $0.864^{s,b\uparrow}$\\
        \cline{2-12}
        & 3ex&- &- & $0.804$& $0.801$& $0.857^{s,b\uparrow}$ &- &- & $0.813^{\star}$& $0.751$& $0.868^{s,b\uparrow}$\\
        \hline
        % \hline
        % & & \multicolumn{10}{c|}{$\Delta\textbf{rater+rubric}$} \\
        % \hline
        % & & \multicolumn{2}{c||}{$\textnormal{LLM}_{H} \to \textnormal{H}_A$}& \multicolumn{3}{c||}{$\textnormal{H}_H \to \textnormal{LLM}_A$} & \multicolumn{2}{c||}{$\textnormal{LLM}_{H} \to \textnormal{H}_A$}& \multicolumn{3}{c|}{$\textnormal{H}_H \to \textnormal{LLM}_A$}\\
        % \hline
        % 1 & & $0.574^{\dagger}$& $0.539$& $0.279$ &$0.271$ & $0.401^{s,b\uparrow}$ & $0.601$& $0.605$& $0.467$ & $0.467$&$0.470$\\
        % \hline
        % 4& & $0.691$& $0.688$& $0.669$ &$0.668$ &$0.682^{s,b\uparrow}$& $0.674$& $0.674$& $0.710^{\star}$ & $0.674$ & $0.707^{b\uparrow}$\\
        % \hline
        % 6 & & $0.645$ & $0.653^{\dagger}$& $0.585$ &$0.595^{\star}$ & $0.644^{s,b\uparrow}$ & $0.660$& $0.670^{\dagger}$& $0.637^{\star}$ &$0.592$ & $0.671^{s,b\uparrow}$\\
        % \hline
    \end{tabular}
    \caption{Kendall's $\tau$ results on AES with GPT-4o and Llama for $\Delta$rater. P. indicates the essay prompt, and C. indicates what ratings are being compared, with ideas, organization, word choice, sentence fluency, and conventions compared for prompt 1, and content, prompt adherence, language, and narrativity compared for prompts 4 and 6. $\tau$ is calculated with singular numerical values for $\Delta$rater and calculated through Pareto dominance comparison for preferences for $\Delta$rubric. Significance tests between separate (sep.), batch (bat.), and edited prompts are performed, where ${^s}$ and ${^b}$ in the edited prompt column represents significant differences with separate and batch prompts respectively. $\dagger$ is indicated next to comparisons that are significantly larger within holistic prompts. $\star$ is indicated next to comparisons that are significantly larger between separate and batch comparisons. $\uparrow$ and $\downarrow$ represent that the $\tau$ value for edited prompts is significantly larger or smaller respectively with the separate $s$ or batch $b$ prompts' $\tau$. The lack of any dagger, star, or arrow denotes no statistical significance. $\textnormal{H}$ is a shortened form for $\textnormal{Human}$. }
    \label{tab:asap_results}
\end{table*}


\begin{table*}[h!]
    \centering
    \begin{tabular}{|c|w{c}{0.3cm}||c|c||c|c|c||c|c||c|c|c|}
    \hline
        & & \multicolumn{5}{c||}{GPT-4o} & \multicolumn{5}{c|}{Llama}\\
        \hline
        \multirow{3}{*}{\rotatebox[origin=c]{90}{$\Delta$rater}} & & \multicolumn{2}{c||}{Holistic} & \multicolumn{3}{c||}{Analytic} & \multicolumn{2}{c||}{Holistic} & \multicolumn{3}{c|}{Analytic}\\
        \cline{3-12}
         & & 0ex & 3ex & Sep. & Bat. & Edited & 0ex & 3ex& Sep. & Bat. & Edited\\
        \cline{3-12}
         & & $0.536$& $0.585$& $0.464^{\star}$& $0.167$&$0.471^{b\uparrow}$ &$0.470$& $0.578$& $0.445^{\star}$& $0.166$&$0.426^{b\uparrow}$\\
        \hline\hline
             \multirow{3}{*}{\rotatebox[origin=c]{90}{$\Delta$rubric}} & Ex. & \multicolumn{2}{c||}{$\textnormal{H}_{H} \to \textnormal{H}_A$}& \multicolumn{3}{c||}{$\textnormal{LLM}_H \to \textnormal{LLM}_A$} & \multicolumn{2}{c||}{$\textnormal{H}_{H} \to \textnormal{H}_A$}& \multicolumn{3}{c|}{$\textnormal{LLM}_H \to \textnormal{LLM}_A$}\\
        \cline{2-12}
        & 0ex & $0.534$&-& $0.640^{\star}$& $0.455$& $0.640^{b\uparrow}$ & $0.551$&- & $0.531^{\star}$& $0.285$& $0.551^{b\uparrow}$\\
        \cline{2-12}
         & 3ex &-&-& $0.623^{\star}$& $0.484$& $0.617 ^{b \uparrow}$ &- &- & $0.566^{\star}$& $0.326$& $0.623^{b\uparrow}$\\
         % \hline
        %  \multirow{2}{*}{\rotatebox[origin=c]{90}{$\Delta$r+r}}& & \multicolumn{2}{|c||}{$\textnormal{LLM}_H \to \textnormal{H}_A$}& \multicolumn{3}{c||}{$\textnormal{H}_H \to \textnormal{LLM}_A$}& \multicolumn{2}{c||}{$\textnormal{LLM}_H \to \textnormal{H}_A$}& \multicolumn{3}{c|}{$\textnormal{H}_H \to \textnormal{LLM}_A$}  \\
        % \cline{3-12}
        % & & $0.541$& $0.530$&  $0.474^{\star}$& $0.34$& $0.459 ^{b\uparrow}$ & $0.478$& $0.545$&  $0.421^{\star}$& $0.183$& $0.424^{b\uparrow}$\\
        
        \hline
    \end{tabular}
    \caption{Kendall's $\tau$ results on IF with GPT-4o and Llama. $\Delta$rubric is calculated through instruction following ratio comparison for preferences. All other calculations and significance follow the methodology of Table~\ref{tab:asap_results}.}
    \label{tab:if_results}
\end{table*}

\subsection{Agreement Calculation}

While direct numerical comparison (e.g., correlation between raw scores) is often calculated through agreement metrics such as Cohen's $\kappa$, autoraters and human have been shown to use the same scales differently, with autoraters often exhibiting compressed or shifted score distributions relative to humans~\cite{kundu2024large}. Preference-based evaluation (e.g., A/B) is more robust to such shifts, as it captures ordinal relationships rather than numerical ones; however, such methods do not scale with large datasets. Addressing both issues, we use Kendall's $\tau$ with tie consideration (calculated using the SciPy \citep{virtanen_scipy_2020} implementation of \citep{df984cd0-5ddb-387f-ae81-1c7b41e92c35}), which operates on the existing numerical scores from the data while evaluating the scores through pairwise ordinal comparisons.
% \jessica{can we actually have a section here that's like: it's "easy" to assign a 1-5 number to a data point if you create a good enough rubric, but that has problems that we said about the distribution shifts of people, but then if you do preference A/B stuff it's "more accurate", but also it's a lot of work because you have to choose what to do preferences over (like you can't just randomly select some points, and you can't do the entire dataset), so this kendall's tau is a middle ground where you are still able to use numbers which is what these datasets had originally, but the evaluation of these numbers now adjusts for distribution shifts (like best of both worlds)}

Each domain employs a different method for aggregating analytic rubric scores when comparing them to holistic rubric scores to calculate concordant pairs for Kendall's $\tau$. In AES, Pareto dominance is used, where essay A is considered better than essay B if all sub-criteria scores of A are at least tied to those of B, with at least one sub-criteria score from A being higher. Pareto dominance is used as a conservative aggregation method to ensure that one essay would truly be better than another without knowing about how the sub-criteria were factored into a holistic rating. In IF, response A is considered to follow instructions better than response B if A has a higher ratio of correctly followed instructions.

Pareto dominance is chosen as a conservative aggregation method precisely because the sub-criteria in ASAP and ASAP++ do not perfectly overlap, as noted in \ref{sec:datasets}. Rather than imposing an arbitrary weighting scheme or score cutoff to determine which essay is better, Pareto dominance avoids assuming any particular trade-off between criteria, thus providing a reliable lower bound on the present agreement.

The confidence interval around the correlations is calculated using bootstrapping \citep{efron1992bootstrap} with 1000 samples. The difference between two conditions is calculated, and the 95\% confidence interval is determined by sampling the 25th and 975th sorted values. For comparisons involving three conditions, the confidence interval is adjusted using Bonferroni correction, with the interval bounds set as the average of the 8th and 9th values and the average of the 991st and 992nd values. The compared conditions are considered significantly different if 0 does not fall within the interval.



% We obtain the holistic and analytic scores from GPT-4 \jessica{cite} with this method. We aim to test three hypotheses regarding prompt formats:
% \begin{enumerate}
%     \item \textbf{Separate} API calls will improve LLM correlation with humans over \textbf{batch} API calls.
%     \item \textbf{Analytic} rubrics will improve LLM correlation with humans over \textbf{holistic} rubrics.
%     \item \textbf{Edited} rubrics (including \textbf{3ex} will improve LLM correlation with humans over all other rubrics because they include best practices for LLM annotations.
% \end{enumerate}

% Additionally, we test several hypotheses that span our conditions:
% \begin{enumerate}
%     \item LLMs will agree more with humans on low-inference decisions than high-inference decisions.
%     \item Rubrics that have more overlap in ASAP with ASAP++ will be closer aligned in preferences.
%     \item Not having examples in certain prompts will cause divergence in alignment over all domains.
%     \item Changing the inference level will cause the least divergence in alignment than changing the decomposition level or generality level.
% \end{enumerate}


% put raw QWK for AES in the appendix

% \begin{enumerate}
%     \item We perform this comparison to determine how much humans and LLMs agree with each other when annotating either holistically or analytically. If this agreement is low, this means that given the same prompt, humans and LLMs do not rate similarly, and this may be an issue with LLM reasoning for the holistic rubric, or that the task is too difficult for the LLM. It may also be the case that LLMs are on a slightly different scale than humans (ex. they are harsher) so an agreement calculation may not be appropriate and instead a correlation would be more appropriate. This is similar to most agreement calculations in previous literature. We calculate agreement between the human ratings and LLM ratings for holistic scoring types. For categorical outputs we perform majority voting over the human annotations if there are multiple annotations and then calculate quadratic weighted kappa (QWK) over the human and LLM annotations, treating the humans as ``one'' human. For any continuous outputs we perform Pearson correlation over the average of the human outputs and the weighted LLM annotations.
%     \item We perform this comparison to determine how much humans (or LLMs) agree between themselves when annotating holistically versus analytically. If this agreement is low, this suggests that there may be a mismatch between the qualities outlined or used by the annotators in the holistic and analytic rubrics. Although the human annotators are different between the holistic annotations and analytic annotations for each domain \jessica{IS THIS TRUE FOR INSTRUCT FOLLOWING? pls confirm}, domain expertise is only a consideration for AES - but the ASAP and ASAP++ annotators seem to have relative equivalence in annotating essays. By calculating agreement  We calculate agreement based on preferences - do the preferences of the holistic scores agree with the preferences of the aggregation of the analytic scores. For AES, we perform Pareto dominance over the analytic scoring as our aggregation method. For emotion recognition we use a KNN trained from \jessica{Tia can you cite} to convert VAD into a categorical emotion to then calculate QWK. For instruct-following we calculate the ratio of followed instructions and calculate a Pearson correlation.
%     \item Lastly, for completion, we perform this comparison to see if there is any correlation between the scoring formats for different scorers (human vs. LLM). Rarely, it may be the case where in the HH/LA case that humans misunderstood the analytic rubric, but the LLM understands the analytic rubric the same as humans interpret the holistic rubric. We expect that if HH/LH, HA/LA, HH/HA, and LH/LA are high, then HH/LA and HA/LH should be high, and if they are low, HH/LA and LA/LH should be low.
% \end{enumerate}

% \begin{enumerate}
%     \item We map both the holistic and analytic prompts given to human annotators such that the prompts are in a format that is palatable for an LLM. 
%     \item We plot the distribution scores of both humans and LLMs - examples of which are shown in the appendix (MAYBE).
%     \item We calculate the comparisons discussed in the next section.
%     \item We plot the distribution scores of humans and LLMs (when not given the data to annotate) to determine what the random scores an LLM would give \jessica{we will not show this in the final paper probably but for curiosity's sake}
%     \item We calculate ICC3K - 2 way mixed effects, consistency metric, mean of k raters \jessica{please check if these apply to your dataset} which looks at ``if raters' scores to the same group of subjects are correlated in an additive manner''. This is similar to Pearson's correlation.
%     \item We calculate a preferences between all pairs of texts within HH/LH and HA/LA to determine even if humans and LLMs do not correlate score-wise if they have the same preferences. For ties between human scores (since they are not continuous) we separate those into another category and also threshold within 0.01 up to 0.5 \jessica{see 12/24 slide} for continuous scores. We also separate the differences between the human scores to see if LLMs agree on preferences with scores between two texts that have larger gaps.
% \end{enumerate}



